\section{A single law defeats arbitrary mixtures}
\label{sec:relu-mixture-extension}

We construct one input law on which every allowed posterior mean has
substantial empirical error against the regression function. The noisy
component has Boolean nuisance coordinates, whereas the signal
component consists of four fixed atoms. The Gaussian-readout sector
then supplies a shared linear witness for all mixtures.

\subsection{The modified law and a gated memorizer}

Fix $p=9/10$, $q=1/10$, an integer $n\ge2$,
$D=\lceil\log_2n\rceil$, $J=2^D$, and $d=D+2$. Thus $n\le J<2n$.
For the exponent statements take $m=\lfloor n^a\rfloor$ and $B=n^b$,
with fixed $a>0$, $b\ge0$, and $Ha+Lb>1/2$. Put
$\delta=2/(\lfloor m/4\rfloor^H B^L)$; sufficiently large $n$
ensures $m\ge4$ and $\delta\le1/2$.
Write $X=(T,A,Z)$ with $Z\in[-1,1]^D$. On the noisy component,
with probability $p$, take $T=0$, $A=1$, $Z$ uniform on
$\{-1,1\}^D$, and $Y$ an independent fair sign. On the signal
component, with probability $q$, take $A=-1$, $Z=0$, $T$ uniform
on $\{-1,-\delta,\delta,1\}$, and $Y=g(T)$, where $g$ is the
bounded tent teacher from Lemma~\ref{lem:tent}. Thus $g(T)$ is again
the exact conditional mean and its risk is $p$. The input covariance is
\[
 \operatorname{diag}\bigl(q(1+\delta^2)/2,4pq,p,\ldots,p\bigr),
\]
so its eigenvalues remain between $1/20$ and one.

\begin{lemma}[A sparse gated memorizer]
\label{lem:relu-gated-memorizer}
Let $r:\{-1,1\}^D\to[-1,1]$ satisfy $\sum_zr(z)^2\le n$.
Set
\[
 M=(D+1)2^D,\quad Q_1=16d(D+1)^2,\quad Q_2=40(D+1)^2n,
 \qquad e=(Q_1Q_2\,6^{L-2})^{1/L}.
\]
At every $N\ge M$ there is an exact network with each weighted layer
energy at most $e$, which equals $r(Z)$ on the noisy atoms and zero
on the signal atoms. It has at most $M$ active first-layer neurons
and at most two active neurons in each later hidden layer.
\end{lemma}
\begin{proof}
For a nonempty set of $j$ Boolean coordinates, let $t$ be their sum,
$t_i=-j+2i$, and $y_i=(-1)^j(-1)^i$. Its parity is represented at
these sums by
\begin{equation}
 2y_0\relu(t+j+1/2)
       -4\sum_{i=0}^{j-1}y_i\relu(t-t_i-1/2).
 \label{eq:relu-margin-parity}
\end{equation}
At $t_0$ its value is $y_0$. On $[t_i,t_i+1/2]$ its slope is
$2y_i$, and on $[t_i+1/2,t_{i+1}]$ its slope is $-2y_i$, so its
next value is $-y_i=y_{i+1}$. This proves the representation by
induction. Each hinge argument has absolute value at least $1/2$ on
the Boolean cube, and the coefficient squared sum is at most
$4+16j\le20(D+1)$. The constant parity is $2\relu(1/2)$.

Apply normalized Fourier expansion and Parseval to $r$. Keeping
separate copies of all parity features gives affine forms $\ell_j(Z)$ and coefficients
$c_j$ satisfying
\[
 r(Z)=\sum_{j=1}^M c_j\relu(\ell_j(Z)),\qquad
 |\ell_j(Z)|\ge1/2,
 \quad T_c:=\sum_jc_j^2\le20(D+1)n/2^D.
\]
Zero-coefficient features fill any unused slots. The constant parity
uses $\ell=1/2$. Replace every affine form by
\[
 \widetilde\ell_j(A,Z)=\ell_j(Z)-2(D+1)(1-A).
\]
At $A=1$ nothing changes. At $A=-1,Z=0$, each form is at most
$D+1/2-4(D+1)<-1/2$, because each original intercept has magnitude
at most $D+1/2$. Thus the features all vanish on the signal atoms.
Every new augmented coefficient norm is at most $4(D+1)$: its squared
norm is at most $D+9(D+1)^2+4(D+1)^2\le16(D+1)^2$.

Partition $N$ into $M$ positive balanced integers $k_j$, so
$k_j\ge N/(2M)$. Use just $M$ active first-layer neurons, with
preactivations $\sqrt{k_j}\,\widetilde\ell_j$.
Their weighted energy is at most $Q_1$.
If $H=1$, give these neurons readout coefficients $c_j/\sqrt{k_j}$.
The output is exactly the desired Fourier sum, and readout energy is
at most $2MT_c\le Q_2$.

If $H\ge2$, let $k_+,k_-$ be positive balanced integers summing to
$N$. The only active second-layer neurons have incoming means
$+\sqrt{k_+}c_j/\sqrt{k_j}$ and
$-\sqrt{k_-}c_j/\sqrt{k_j}$. Put
$F=\sum_jc_j\relu(\widetilde\ell_j)$.
Their activations are $\sqrt{k_+}\relu(F)$ and
$\sqrt{k_-}\relu(-F)$, and the second-layer energy is at most
$2MT_c\le Q_2$.
For each later hidden layer, use only two active rows, chosen so that
its normalized preactivations are $\pm(h_+/\sqrt{k_+}-h_-/\sqrt{k_-})$.
Explicitly, its $2$-by-$2$ nonzero block is
\[
 \begin{pmatrix}
 1&-\sqrt{k_+/k_-}\\
 -\sqrt{k_-/k_+}&1
 \end{pmatrix}.
\]
Its squared Frobenius norm is
$N(1/k_++1/k_-)\le6$. The normalized activation difference remains
exactly $F$, by $\relu(t)-\relu(-t)=t$.
Readout weights $1/\sqrt{k_+}$ and $-1/\sqrt{k_-}$ recover $F$,
and their weighted energy is at most six.

The product of the $L$ weighted layer energies is at most
$Q_1Q_2\,6^{L-2}$. Positive homogeneity balances them to their
geometric mean without altering the output or active supports. If
one is zero, the target function is zero and the zero network suffices.
This proves all assertions.
\end{proof}

Apply Lemma~\ref{lem:relu-profile-comparison} to this sparse
memorizer. Its covariance cost now has coefficient
\[
 D_*:=d+2+2(H-1),
\]
independent of $M$. This is why no enlarged width threshold is
needed to include mean-field scaling.

\subsection{One sample event and empirical separation}

Keep $\tau_n=n^{-1/L}d^{-4}$ and set
\[
 N_*(n)=\max\{m,(D+1)2^D\}=O(m+n\log n).
\]
For every sufficiently large $n$, the following properties hold on
one sample event, uniformly over $N\ge N_*$,
$0<\sigma^2\le1/N$, and $\lambda>0$.

First, the singleton occupancy and signal-count bounds remain
\[
 S_1/n\ge pe^{-p}-n^{-1/3},\qquad
 n_s/n\le q+n^{-1/3},
\]
with failure probability at most $2e^{-n^{1/3}/2}$. Indeed
$\mathbb E S_1=np(1-p/J)^{n-1}\ge npe^{-p}$ because $J\ge n$;
changing one observation changes $S_1$ by at most two. McDiarmid and
Hoeffding give the two bounds.
For $\lambda\le\tau_n$, use the empirical noisy means as $r$ in
Lemma~\ref{lem:relu-gated-memorizer}. The memorizer has energy at
most $Le=O_H(n^{1/L}d^{5/L})$ and is bounded by one on the support.
Set $K_e=\sqrt2 L(e+1)^{H/2}$ and
$r_n=[8n(K_e+1)]^{-1}$ in the selective Gaussian comparison lemma.
Its expected empirical loss is at most the exact memorizer's loss
plus $1/n$. The Gibbs variational identity gives
\[
 \int\widehat R_S(z_w)d\Pi
 \le\widehat R_S(z_{w_*})+1/n
       +\lambda\{Le+2D_*\log(1/r_n)\}.
\]
Here $\log(1/r_n)=O_H(\log n)$, uniformly in learner width and
prior scale. After subtracting the within-atom noise residual,
the expected weighted noisy residual is therefore at most
$q+o(1)\le1/5$.
Jensen and the reverse triangle inequality imply that the raw
posterior mean $f$ obeys
\[
 \frac1n\sum_{i\text{ noisy}}f(X_i)^2
 \ge\left(\sqrt{1/3}-\sqrt{1/5}\right)^2>1/160.
\]
The same conclusion holds after clipping this mean to $[-1,1]$,
because clipping cannot increase its noisy-label fitting residual.

For $\lambda\ge\tau_n$, normalize the matrices by
$A_1=\sqrt{d/N}\,W^1$, $A_\ell=W^\ell$ internally, and
$A_L=\sqrt N\,v^{\mathsf T}$. Every coordinate then has variance
$\sigma^2$. Put $T=\max_\ell\|A_\ell\|_{\rm op}$ and
$M_d=1+\sqrt{d+1}$. Gaussian operator-norm concentration and
$N\sigma^2\le1$ give
$\mu\{T>M_d+t\}\le L e^{-t^2/(2\sigma^2)}$.
The prior expected empirical loss is at most three, so
$Z\ge e^{-3/(2\lambda\sigma^2)}$ by Jensen. Dividing the prior
tail by $Z$ and integrating yields
\[
 \left(\int T^L d\Pi\right)^{1/L}
 \le M_d+\sqrt{3/\lambda+2\log L}+2\sqrt L.
\]
The input-coordinate Lipschitz constant is at most $d^{-1/2}T^L$.
The inner signal points have $A=-1,Z=0$, and therefore
\[
 \frac{|f(\delta,-1,0)-f(-\delta,-1,0)|}{2}=o(1)\le1/2
\]
whenever $Ha+Lb>1/2$. Clipping the raw mean preserves this bound.
Require in addition that each inner signal atom occurs at least
$qn/8$ times. Its count has mean $qn/4$, so the additional failure
probability is at most $2e^{-nq/32}=2e^{-n/320}$.
The sum of the two squared errors against their labels $\pm1$ is
at least $1/2$. Therefore these observations alone give empirical
squared error against $g$ at least $qn/16=n/160$.

Thus every allowed raw posterior mean, and every clipped raw
posterior mean, has empirical squared error against $g$ at least
$1/160$ on this event. This conclusion, unlike an individual
population-risk lower bound, can be combined with a common linear
separator.

\begin{theorem}[Failure of all data-dependent mixtures]
\label{thm:relu-mixture-failure}
Fix $H\ge1$, $a>0$, $b\ge0$, and $Ha+(H+1)b>1/2$.
For all sufficiently large integers $n$, the modified law above has a
conditional mean of width at most $\lfloor n^a\rfloor$, depth
$L=H+1$, coefficient bound $n^b$, and output magnitude at most one.
With sample probability at least
\[
 1-2e^{-n^{1/3}/2}-2e^{-n/320}-1/n,
\]
every data-dependent probability mixture of ordinary raw posterior
means over
\[
 N\ge N_*(n),\qquad 0<\sigma^2\le1/N,\qquad \lambda>0
\]
has excess population squared risk at least $1/300000$.
The same is true if each component mean is clipped to $[-1,1]$
before mixing, or if the mixture contains both types of components.

This includes every canonical exponent $\gamma\ge1$ at
$\sigma^2=N^{-\gamma}$, with no limiting operation. The same
conclusion holds for mixtures of strong-prior global minimizers.
\end{theorem}
\begin{proof}
The previous subsection proves the empirical separation. The
Gaussian-readout sector and separator lemmas apply because
$Y_i g(X_i)=g(X_i)^2$ coordinatewise: $g$ is zero on noisy atoms and
is exactly the label on signal atoms. Add the event
$\sum_j k_j^2\le4n$ for noisy counts. It has failure probability at
most $1/n$: if $C=\sum_j\binom{k_j}{2}$ counts noisy collisions,
then $\sum_jk_j^2\le n+2C$, $\mathbb EC\le n/2$, and
\[
 \operatorname{Var}(C)
 \le n\{p^2/2+p^3(1-p)\}<n.
\]
Indeed disjoint collision indicators are independent, and pairs
sharing one observation have covariance $p^3(1-p)/J^2$. Since $J\ge n$, the displayed variance bound follows.
Chebyshev now proves the claim.
For the separator $h$, the noisy contribution to its squared
population norm is at most
$J\sum_j k_j^2/(n^2p)\le8/p$. The signal contribution is at most
$4n_s^2/(n^2q)\le16q$, since $n_s\le2qn$ eventually.
Thus $\|h\|_2^2\le8/p+16q<11$. Applying the separator with $c=1/160$ gives excess at
least $1/(25600\cdot11)>1/300000$, uniformly over all mixtures.

For strong-prior global minimizers, comparison with the deterministic
sparse memorizer gives the same low-temperature residual estimate.
Comparison with zero gives $Q_N\le1/\lambda$ at high temperature,
and hence the same vanishing inner-pair half-difference. Every
minimizing readout is a ridge solution and satisfies the sector
inequality. Thus these components can be included in the same
mixture family. All ingredients use the same sample event.
\end{proof}

The mixture may be chosen using the sample, validation data, or the
unknown law itself. In particular, a hyperprior over the allowed
likelihood temperatures does not repair this counterexample. The
claim concerns mixtures of raw posterior means and clipping of each
such mean before mixing. It does not assert the same property for
averaging clipped posterior draws, or for clipping only after taking
an arbitrary mixture of raw means.


\begin{corollary}[Any fixed positive localized noise mass]
Fix $0<p\le9/10$. The same strict comparator-growth condition admits
one law of Bayes risk $p$ for which every mixture in
Theorem~\ref{thm:relu-mixture-failure} has excess risk at least
\[
 \frac{p^3e^{-2p}}{600000}.
\]
The width threshold remains $O(m+n\log n)$, and the sample probability
is at least
$1-2e^{-n^{1/3}/2}-2e^{-nr/16}-1/n$, where $r=pe^{-p}/32$.
The sufficiently large sample threshold may depend on $p$.
\end{corollary}
\begin{proof}
Keep noisy mass $p$ with Boolean $Z$. Give the inner pair total mass
$r=pe^{-p}/32$, and the outer anchors total mass $1-p-r$; all signal
atoms have $A=-1,Z=0$. The same teacher is the conditional mean.
The covariance is diagonal with entries
$1-p-r+r\delta^2$, $4p(1-p)$, and $p$ in the nuisance coordinates,
so it stays nondegenerate for each fixed $p$.

Use the singleton event together with Hoeffding control of the
inner-pair count, whose sample fraction is at most $r+n^{-1/3}$.
These two events have total failure at most $2e^{-n^{1/3}/2}$.
Write $s=pe^{-p}$. The singleton count is at least $sn/2$ eventually,
and the memorizer's empirical signal loss is only the inner-pair
count. The low-temperature noisy fitting residual is therefore at
most $s/16$ for sufficiently large $n$. Its noisy prediction norm
squared, divided by $n$, is at least
$s(1/\sqrt2-1/4)^2>s/256$.
Each inner signal atom has count at least $rn/4$, except with total
probability $2e^{-nr/16}$. The high-temperature half-difference is
at most $1/2$, so its empirical squared error is at least $rn/8$,
which equals $sn/256$.
Thus the common separator applies with $c=s/256$.
Its noisy squared norm is at most $8/p$. Since the inner-pair count
is at most $2rn$ eventually, its signal squared norm is at most
$8r$. Hence $V\le8/p+8r\le9/p$, and
\[
 c^2/V\ge p^3e^{-2p}/589824>p^3e^{-2p}/600000.
\]
The occupancy and collision arguments are unchanged. Labels retain
unit magnitude on the noisy region and inner pair; this is localized
regression noise, not independent Boolean label flips.
\end{proof}
